% GENERATED FILE. Edit canonical lessons, then run npm run generate:books.
% canonical-source-hash: fnv1a64:9933170f77f0fd9f
% canonical-lessons: FR-C332-passoire, FR-C332-egoutter, FR-C332-saupoudrer, FR-C332-assaisonner, FR-C332-deguster

\chapter{Assaisonner}
\label{ch:fr-assaisonner}
\hlchaptermodality{332}

\begin{chapteropening}
\textbf{By the end of this chapter:} \emph{I can describe finishing a dish: the colander, draining, sprinkling, seasoning and tasting.}

Complete the last lesson of chapter 332: I can describe finishing a dish: the colander, draining, sprinkling, seasoning and tasting.
\end{chapteropening}

\section[la passoire]{la passoire — the colander, the sieve}

\label{lesson:FR-C332-passoire}



\begin{quote}
Before the new one: say the French for to turn golden, then the French for to brown on top.
\end{quote}

\subsection*{You'll want to know: la passoire}
\textbf{la passoire} — \textquotedblleft{}the colander, the sieve\textquotedblright{}.

\begin{quote}
Il met les pâtes dans la passoire.
\end{quote}

\begin{quote}
\emph{He puts the pasta in the colander.}
\end{quote}

Informal: \emph{une mémoire comme une passoire} — a memory like a sieve.

\begin{cousinweb}
\emph{passoire} is built on \emph{passer} \textquotedblleft{}to pass\textquotedblright{}, which is inherited from Vulgar Latin \emph{passāre}, built on Latin \emph{passus} \textquotedblleft{}a step\textquotedblright{} (Lewis \& Short).
\end{cousinweb}

\begin{grammarlens}[title={What you've built}]
One more word for telling, explaining and describing at B1.
\end{grammarlens}

\subsection*{Your turn}
\begin{itemize}
\raggedright
  \item \emph{Say it:} \emph{la passoire}
  \item \emph{Say it:} \emph{la passoire}, once more
  \item \emph{Say it:} a sentence of your own with \emph{la passoire}
  \item \emph{Recall:} say the French for to turn golden, then the French for to brown on top, then say \emph{la passoire} again
\end{itemize}

\begin{tcolorbox}[breakable,colback=teal!4,colframe=teal!35!black,title={Before you move on}]
What does la passoire mean? (\textquotedblleft{}The colander, the sieve\textquotedblright{}.) Say it once more.
\end{tcolorbox}
\section[égoutter]{égoutter — to drain}

\label{lesson:FR-C332-egoutter}



\begin{quote}
Before the new one: say the French for to brown on top, then the French for the colander.
\end{quote}

\subsection*{You'll want to know: égoutter}
\textbf{égoutter} — \textquotedblleft{}to drain\textquotedblright{}.

\begin{quote}
Elle égoutte les pâtes.
\end{quote}

\begin{quote}
\emph{She drains the pasta.}
\end{quote}

\emph{Laisser égoutter} — to leave to drain.

\begin{cousinweb}
\emph{égoutter} is built on \emph{goutte} \textquotedblleft{}a drop\textquotedblright{}, which is inherited from Latin \emph{gutta} \textquotedblleft{}a drop\textquotedblright{} (Lewis \& Short).
\end{cousinweb}

\begin{grammarlens}[title={What you've built}]
One more word for telling, explaining and describing at B1.
\end{grammarlens}

\subsection*{Your turn}
\begin{itemize}
\raggedright
  \item \emph{Say it:} \emph{égoutter}
  \item \emph{Say it:} \emph{égoutter}, once more
  \item \emph{Say it:} a sentence of your own with \emph{égoutter}
  \item \emph{Recall:} say the French for to brown on top, then the French for the colander, then say \emph{égoutter} again
\end{itemize}

\begin{tcolorbox}[breakable,colback=teal!4,colframe=teal!35!black,title={Before you move on}]
What does égoutter mean? (\textquotedblleft{}To drain\textquotedblright{}.) Say it once more.
\end{tcolorbox}
\section[saupoudrer]{saupoudrer — to sprinkle, to dust}

\label{lesson:FR-C332-saupoudrer}



\begin{quote}
Before the new one: say the French for the colander, then the French for to drain.
\end{quote}

\subsection*{You'll want to know: saupoudrer}
\textbf{saupoudrer} — \textquotedblleft{}to sprinkle, to dust\textquotedblright{}.

\begin{quote}
Il saupoudre le gâteau de sucre.
\end{quote}

\begin{quote}
\emph{He dusts the cake with sugar.}
\end{quote}

\emph{Saupoudrer de} — to dust with.

\begin{cousinweb}
\emph{saupoudrer} is built from \emph{sel} \textquotedblleft{}salt\textquotedblright{} and \emph{poudrer} \textquotedblleft{}to powder\textquotedblright{} (EtymDB); \emph{poudre} goes back to Latin \emph{pulvis} \textquotedblleft{}dust\textquotedblright{}.
\end{cousinweb}

\begin{grammarlens}[title={What you've built}]
One more word for telling, explaining and describing at B1.
\end{grammarlens}

\subsection*{Your turn}
\begin{itemize}
\raggedright
  \item \emph{Say it:} \emph{saupoudrer}
  \item \emph{Say it:} \emph{saupoudrer}, once more
  \item \emph{Say it:} a sentence of your own with \emph{saupoudrer}
  \item \emph{Recall:} say the French for the colander, then the French for to drain, then say \emph{saupoudrer} again
\end{itemize}

\begin{tcolorbox}[breakable,colback=teal!4,colframe=teal!35!black,title={Before you move on}]
What does saupoudrer mean? (\textquotedblleft{}To sprinkle, to dust\textquotedblright{}.) Say it once more.
\end{tcolorbox}
\section[assaisonner]{assaisonner — to season}

\label{lesson:FR-C332-assaisonner}



\begin{quote}
Before the new one: say the French for to drain, then the French for to sprinkle.
\end{quote}

\subsection*{You'll want to know: assaisonner}
\textbf{assaisonner} — \textquotedblleft{}to season\textquotedblright{}.

\begin{quote}
Elle assaisonne la salade avec de l'huile et du vinaigre.
\end{quote}

\begin{quote}
\emph{She dresses the salad with oil and vinegar.}
\end{quote}

\emph{L'assaisonnement} — the seasoning, the dressing.

\begin{cousinweb}
\emph{assaisonner} is built on \emph{saison} \textquotedblleft{}season, the right time\textquotedblright{}, usually traced to Latin \emph{satiō} \textquotedblleft{}a sowing\textquotedblright{} (Lewis \& Short): to bring food to the right point.
\end{cousinweb}

\begin{grammarlens}[title={What you've built}]
One more word for telling, explaining and describing at B1.
\end{grammarlens}

\subsection*{Your turn}
\begin{itemize}
\raggedright
  \item \emph{Say it:} \emph{assaisonner}
  \item \emph{Say it:} \emph{assaisonner}, once more
  \item \emph{Say it:} a sentence of your own with \emph{assaisonner}
  \item \emph{Recall:} say the French for to drain, then the French for to sprinkle, then say \emph{assaisonner} again
\end{itemize}

\begin{tcolorbox}[breakable,colback=teal!4,colframe=teal!35!black,title={Before you move on}]
What does assaisonner mean? (\textquotedblleft{}To season\textquotedblright{}.) Say it once more.
\end{tcolorbox}
\section[déguster]{déguster — to taste, to savour}

\label{lesson:FR-C332-deguster}



\begin{quote}
Before the new one: say the French for to sprinkle, then the French for to season.
\end{quote}

\subsection*{You'll want to know: déguster}
\textbf{déguster} — \textquotedblleft{}to taste, to savour\textquotedblright{}.

\begin{quote}
Nous avons dégusté les fromages de la région.
\end{quote}

\begin{quote}
\emph{We tasted the cheeses of the region.}
\end{quote}

To eat or drink slowly, with pleasure. Informal: \emph{il a dégusté} — he really suffered.

\begin{cousinweb}
\emph{déguster} is a learned borrowing from Latin \emph{dēgustāre} \textquotedblleft{}to taste\textquotedblright{} (EtymDB), built on \emph{gustus} \textquotedblleft{}taste\textquotedblright{}.
\end{cousinweb}

\begin{grammarlens}[title={What you've built}]
One more word for telling, explaining and describing at B1.
\end{grammarlens}

\subsection*{Your turn}
\begin{itemize}
\raggedright
  \item \emph{Say it:} \emph{déguster}
  \item \emph{Say it:} \emph{déguster}, once more
  \item \emph{Say it:} a sentence of your own with \emph{déguster}
  \item \emph{Recall:} say the French for to sprinkle, then the French for to season, then say \emph{déguster} again
\end{itemize}

\begin{tcolorbox}[breakable,colback=teal!4,colframe=teal!35!black,title={Before you move on}]
What does déguster mean? (\textquotedblleft{}To taste, to savour\textquotedblright{}.) Say it once more.
\end{tcolorbox}
